\section{The actual entry space and its flat local deformation charts}
\label{sec:entry-space-flat}

Fix the arithmetic frame of
Section~\ref{sec:arithmetic-frame-uniform-depth},
with \(\operatorname{Frame.Ok}(C_0)\) and
\(\operatorname{Frame.Big}(C_0)\).
Write \(d=\sum_vd_v\), let \(\mathcal O\) be a complete
coefficient DVR with residue field \(k\), and use ranks
\(n=2,3\).
The weights are \(w=(1,-1)\) and \(w=(1,0,-1)\),
respectively; \(J\) is the reversal matrix, and
\(g^*=cg^{-1}c^{-1}\).

The normalized representations used here consist of
actual matrices and ordered local characters satisfying
\[
 \begin{split}
 &\rho(gh)=\rho(g)\rho(h),\qquad
                 \rho(g^*)=J\,{}^t\rho(g)J,\\
 &(\rho(g_0)-\rho(g_0^*))_{ij}=0\quad(i\ne j),\\
 &\prod_{i=0}^{n-1}
       \bigl(\rho(\operatorname{loc}_v(g_i))
                         -\psi_{vi}(g_i)I_n\bigr)=0
                                  \quad\hbox{for every }(g_i).
 \end{split}                                                  \tag{LS1}
\]
The product is in the displayed order.
The residual diagonal and local characters are
\(\Psi^{w_i}\).
The ideals \(\mathfrak t_+,\mathfrak t_-\) are generated
by entries below and above the diagonal.
In rank three \(\mathfrak e\) is generated by the four
outer--middle entries.
Type I means off both triangular loci and off
\(V(\mathfrak e)\); Type E means off both triangular
loci and on \(V(\mathfrak e)\).

\subsection{Construction and algebraization of the one entry space}
There is a complete Noetherian local ring
\(R^{\rm diag}\) representing continuous data (LS1)
with the diagonal residual representation.
Here is the actual construction.
For the continuous framed representation functor
with its fixed diagonal residue, the condition
\(\Phi_\ell\) makes the relative tangent space finite.
The functor preserves fibre products, including the
fibre products used in the infinitesimal uniqueness
condition: a pair of compatible matrices has a
unique matrix lift to the fibre product.
The standard successive small-extension construction
therefore gives its universal complete local ring.
There are no automorphisms to identify, since the
basis is part of the data.
Lifting a tangent basis gives a surjection from a
finite-variable \(\mathcal O\)-power-series ring.

Adjoin the independent ordered local characters.
For each \(v\), their deformation ring is the
completed group algebra of the maximal pro-\(\ell\)
quotient of \(G_v^{\rm ab}\), translated by their
fixed Teichm\"uller residue characters.
Local class field theory makes that quotient a
finite \(\ell\)-group times \(\mathbb Z_\ell^{d_v+1}\).
Thus the completed tensor product of these finitely
many character rings with the framed ring is again
complete Noetherian.
In this ring take the ideal generated by every
entry of the polarization, marking, and ordered
product equations in (LS1).
This ideal is closed: every ideal of a complete
Noetherian local ring is closed.
Its quotient \(R^{\rm diag}\) represents precisely
the continuous data (LS1).
Indeed a continuous map kills the ideal exactly
when it kills all its displayed generators, and
all equations are preserved under the inverse
limit of the finite coefficient quotients.

For each ordinary place choose \(s_v\) with
\(\Psi(\operatorname{loc}_v s_v)^2\ne1\).
Hensel diagonalization at its distinct residual
eigenvalues gives inverse matrices \(P_v,Q_v\),
congruent to one, diagonalizing \(\rho(s_v)\).
Impose the polarization and marked-element
equations and the entries below the diagonal of
\(Q_v\rho(\operatorname{loc}_v s)P_v\), for all \(s\).
Its diagonal entries give the ordered characters.
Conversely (LS1) implies these same equations.
Set every \(g_i=s_v\) except one.
The other factors are diagonal.
The \((a,b)\)-entry for \(b\le a\) isolates
\[
       \bigl(Q_v\rho(\operatorname{loc}_v s)P_v
                            -\psi_{va}(s)I_n\bigr)_{ab}
\]
times products of unit eigenvalue differences.
For \(n\le3\), direct multiplication of the at
most three factors proves that it is zero.
Thus this quotient represents the stated ordinary
functor, including its ordered characters.

Give \(\rho(g)_{ij}\) weight \(w_i-w_j\) and each
local character weight zero.
Conjugation by \(\operatorname{diag}(t^{w_i})\)
preserves (LS1) and gives a completed torus grading.
Define the actual source rings
\[
 \begin{split}
 B&=\overline{\mathcal O[
       \text{weight-zero entry monomials},\,\psi_{vi}(s)]}
                                \subset R^{\rm diag},\\
 A&=B[\rho(g)_{ij}:g,i,j]\subset R^{\rm diag}.
 \end{split}                                                  \tag{LS2}
\]
The closure is in the maximal-ideal topology of
\(R^{\rm diag}\).

We record their finiteness and topology.
The entries and characters topologically generate
\(R^{\rm diag}\), by the representing property.
Choose finitely many homogeneous such elements
spanning its relative cotangent space.
They give a graded surjection from
\(\mathcal O[[X_1,\ldots,X_h]]\).
The exponent semigroup of weight-zero monomials
is finitely generated.
For each fixed weight the exponent vectors
have finitely many coordinatewise minimal
elements, by Dickson's lemma; every other
exponent is one of these plus a weight-zero
exponent.
Decomposing a formal series accordingly proves
that the degree-zero part is complete Noetherian
and each fixed-weight part is finite over it.
These assertions pass to the homogeneous quotient,
by projecting coefficients to their weight.
Thus the degree-zero part is exactly \(B\),
each \(A_j\) is finite over \(B\), and \(A\) is
a finitely generated graded \(B\)-algebra.
In particular \(A\) is Noetherian.
On each \(A_j\) the induced topology is cofinal
with the \(\mathfrak m_B\)-adic topology: a long
fixed-weight monomial contains arbitrarily
many nonconstant weight-zero factors.

This one algebra also represents the points
with nondiagonal triangular residue.
For clarity we give the actual algebraization
argument.
Suppose \(Q\) is complete local Noetherian with
finite residue extension of \(k\), and \(\rho_Q\)
satisfies (LS1), continuously, with residual
diagonal entries and local characters
\(\Psi^{w_i}\).
Its residual representation is upper or lower
triangular.
Here are the coordinate equations for this assertion.
Write \(a_{ij}(g)=\bar\rho_Q(g)_{ij}\).
Each diagonal function is the character
\(\Psi^{w_i}\), so for every \(g,h\)
\[
 \sum_{j\ne i}a_{ij}(g)a_{ji}(h)=0.
\]
The marked matrix
\[
 X=\bar\rho_Q(g_0)-\bar\rho_Q(g_0^*)
                         =\operatorname{diag}(\zeta,0,-\zeta)
\]
has \(\zeta\ne0\).
Comparing the \(i\)-th diagonal entries of
\(\bar\rho_Q(g)X\bar\rho_Q(h)\) gives the second equation
\[
 \sum_{j\ne i}X_{jj}a_{ij}(g)a_{ji}(h)=0.
\]
Since \(2\zeta\ne0\), these two equations imply
\[
 a_{01}(g)a_{10}(h)=
 a_{02}(g)a_{20}(h)=
 a_{12}(g)a_{21}(h)=0
                              \quad\text{for all }g,h.
\]
Polarization gives
\(a_{01}(g^*)=a_{12}(g)\) and
\(a_{10}(g^*)=a_{21}(g)\).
Consequently either \(a_{01}=a_{12}=0\) identically,
or \(a_{10}=a_{21}=0\) identically.
Also either \(a_{02}=0\) or \(a_{20}=0\) identically.

Only the two mixed cases require another equation.
If \(a_{01}=a_{12}=a_{20}=0\), the \((0,1)\)-entry
of \(\bar\rho_Q(gh)=\bar\rho_Q(g)\bar\rho_Q(h)\) gives
\(a_{02}(g)a_{21}(h)=0\).
Either \(a_{02}=0\), giving lower triangularity,
or \(a_{21}=a_{10}=0\), giving upper triangularity.
If \(a_{10}=a_{21}=a_{02}=0\), the \((2,1)\)-entry
instead gives \(a_{20}(g)a_{01}(h)=0\);
its two alternatives give upper or lower
triangularity in the same way.
The other cases are already triangular.
In rank two the one diagonal equation is
\(a_{01}(g)a_{10}(h)=0\), which proves the assertion.
For points of the fibre in (LS2), the same
conclusion follows immediately: nonzero entries
of opposite weights would give a nonzero
weight-zero product in the residual augmentation.

Put \(Q_0=\{q\in Q:\bar q\in k\}\).
It is complete local with residue \(k\).
The ring \(Q\) is finite over \(Q_0\):
lift a residue basis including \(1\); subtract
its residue combination from any element.
The remainder belongs to
\(\mathfrak m_Q\subset Q_0\).
These lifts generate \(Q\) over \(Q_0\).
Noetherianity descends along this finite extension,
and the maximal ideals of \(Q_0,Q\) have the
same underlying set and the same powers.

Let \(P_j\subset Q\) be the \(Q_0\)-span of
entry monomials of weight \(j\), and form
\[
                          D=\bigoplus_jP_jt^j
                                    \subset Q[t,t^{-1}].      \tag{LS3}
\]
Every weight-zero off-diagonal monomial belongs
to \(\mathfrak m_Q\subset Q_0\); diagonal and
local-character entries belong to \(Q_0\).
Thus \(P_0=Q_0\), and the polynomials whose
degree-zero coefficient lies in
\(\mathfrak m_{Q_0}\) form a maximal ideal
\(\mathfrak n\) of \(D\).
Products of opposite nonzero degrees have that
property.
Each entry span is finite over \(Q_0\), so
\(D\) is a finitely generated \(Q_0\)-algebra.
Each \(P_j\) is finite over \(Q_0\) by the same
fixed-weight monomial decomposition.
Its induced \(\mathfrak n\)-adic topology is
cofinal with its
\(\mathfrak m_{Q_0}\)-adic topology.
Indeed remove primitive zero-weight factors
from a long fixed-weight monomial.
There are finitely many such factors, of
bounded length; each lies in
\(\mathfrak m_{Q_0}\), and the remaining
minimal monomial has bounded length.

In \(\widehat D_{\mathfrak n}\), the entries
\(t^{w_i-w_j}\rho_Q(g)_{ij}\), together with
the degree-zero characters, give continuous
residually diagonal data (LS1).
The universal map from \(R^{\rm diag}\) exists.
Its image on \(B\) lies in the closed
degree-zero part \(Q_0\), and its image on
\(A\) lies in (LS3).
Evaluation at \(t=1\) gives
\[
                              A\longrightarrow Q              \tag{LS4}
\]
with exactly the original entries and characters,
and with \(\mathfrak m_B\) mapping into
\(\mathfrak m_Q\).
Continuity on each entry lattice follows from
Artin--Rees inside the finite \(Q_0\)-module
\(Q\).
Uniqueness holds on entries and characters,
then on \(B\) by fixed-weight adic approximation,
then on \(A\).

For locally constant lower data over an arbitrary
\(k\)-field or over \(k[T]\), their entries have
finite image and hence finite \(k\)-span.
All nonconstant weight-zero monomials vanish.
Use \(Q_0=k\) in (LS3).
The same vertex completion and evaluation give
the actual field and polynomial points of \(A\).
Consequently the previously constructed lower
interpolation families are families in this
one entry space.
Their connected core and dense nonexceptional
locus require no further point-identification
hypothesis.

\subsection{Flat local charts and their dimension formula}
Let \(x\) be a closed point of \(A\) above
\(\mathfrak m_B\), with residue field \(k'\).
Choose the finite unramified coefficient extension
\(\mathcal O'/\mathcal O\) with residue field \(k'\).
For its specified residue map let
\[
 R_x=\widehat{(A\otimes_{\mathcal O}\mathcal O')_{\mathfrak n_x}}.
                                                               \tag{LS5}
\]
The map \(A\to R_x\) is flat: coefficient extension,
localization, and Noetherian local completion
are flat.
After localizing \(A\) at \(x\), it is faithfully
flat.
By (LS4) and its uniqueness, (LS5) represents
all normalized deformations of the actual
residue representation at \(x\).
This proves the source's \srcid{DefIso} statement;
it does not assume a new unrelated deformation
ring.

Minimal primes of \(R_x\) map to generic points
of components through \(x\), and every such
component has a minimal-prime lift.
Going down and faithful flatness prove these
two assertions.
Thus (LS5) gives the required flat chart.

For a prime \(P\subset R_x\), put \(p=P\cap A\),
\(p_B=p\cap B\) and
\(t=\operatorname{trdeg}_{B/p_B}(A/p)\).
Then
\[
                       \dim R_x/P\le\dim B/p_B+t.             \tag{LS6}
\]
For each \(p\subseteq x\), equality is attained
at some prime of \(R_x\) above \(p\).
The finite coefficient extension preserves
transcendence degree.
The finite-type dimension formula at the closed
point gives
\(\dim(A_x/pA_x)=\dim B/p_B+t\).
Completion preserves this quotient dimension;
excellence makes the completion of the local
domain \(A_x/pA_x\) equidimensional.
Any prime over it has at most that dimension,
and its minimal primes have exactly that
dimension.
This proves (LS6) with its existence assertion.

\subsection{The actual tangent and obstruction maps}
Write \(R=R_x\).
Extend its residue representation to the unitary
semidirect product on the real group, with
complex conjugation of matrix part \(I_n\).
The adjoint action on the CM subgroup is
conjugation by \(\bar\rho\); on the other coset
it includes \(\vartheta(X)=-J\,{}^tXJ\).
Choose the actual ordinary flags and denote
their Borel Lie algebras by \(b_v\).
All cone groups below include local correcting
vectors as in (AF9).

For a small extension \(C'\twoheadrightarrow C\)
with kernel \(I\) and \(\mathfrak m_{C'}I=0\),
choose continuous setwise global matrix lifts
\(\widetilde r(s)\), flag lifts
\(\widetilde g_v\), and upper-triangular
local lifts \(\widetilde t_v(s)\).
They exist by finite sections on coefficient
quotients.
The defects are the matrices with entries in \(I\)
\[
 \begin{split}
 \Omega(s,t)&=\widetilde r(s)
    \operatorname{sgnAct}_{\delta(s)}(\widetilde r(t))
                                  \widetilde r(st)^{-1}-I_n,\\
 a_v(s)&=\widetilde r(\operatorname{loc}_v s)
   \bigl(\widetilde g_v^{-1}\widetilde t_v(s)
                                      \widetilde g_v\bigr)^{-1}
                                                   -I_n .
 \end{split}                                                  \tag{LS7}
\]
Associativity gives \(d\Omega=0\), and local
multiplication gives
\(\Omega|_{G_v}-da_v\in b_v\).
Applying \(u:I\to k'\) gives a class in
\(H_C^2(M_n,b)\).

If this class is zero, write \(u\Omega=d\beta\)
and
\(ua_v-\beta|_{G_v}
       +(\operatorname{Ad}(\bar r)-1)Y_v\in b_v\).
In the one-dimensional pushout selected by \(u\),
replace the global lifts by
\((1-\beta(s))\widetilde r(s)\), and the flags by
\(\widetilde g_v(1-Y_v)\).
Square-zero multiplication proves the corrected
representation and the ordinary condition.
For complex conjugation normalize \(1+X\),
where \(J\,{}^tXJ=X\), by \(1-X/2\).
For the marked difference use
\(Y_{ij}=-E_{ij}/(\lambda_j-\lambda_i)\)
on its off-diagonal error.
The denominators are units and its
antidiagonal symmetry makes \(\vartheta(Y)=Y\);
this conjugation preserves the first normalization.
Thus zero obstruction gives an actual
normalized ordinary lift.

Choose a minimal relative presentation
\(P=\mathcal O'[[X_1,\ldots,X_t]]\twoheadrightarrow R\),
with kernel \(I_R\).
The obstruction gives an injection
\[
              (I_R/\mathfrak m_PI_R)^*
                   \lhook\joinrel\longrightarrow H_C^2(M_n,b).
                                                               \tag{LS8}
\]
Indeed a nonzero kernel functional would give
a nonzero one-dimensional quotient of
\(P/\mathfrak m_PI_R\to R\) with a section.
A section is impossible for a minimal
presentation.
For a relation write
\(f=f(0)+\sum_iX_i h_i\), with
\(h_i\in\mathfrak m_P\).
Subtract its value under the section.
Every kernel element is then in the maximal
ideal times that kernel; Nakayama kills it.
This also treats a coefficient-uniformizer
relation.

A cone \(H^1\)-class gives first-order matrices
\((1+\epsilon z(s))\bar r(s)\) with their
corrected local flags.
Normalize them as above.
This shows that every class comes from a
relative tangent vector.
The torus matrix \(T_0=\operatorname{diag}(w_i)\)
is nonzero and its first-order conjugation
preserves both normalizations.
If its cocycle is nonzero it is a nonzero
tangent direction killed in \(H_C^1\);
if zero it is a nonzero cone invariant.
Consequently
\[
                    h_C^1+1\le t+h_C^0 .                     \tag{LS9}
\]
The local boundary rigidity in this comparison
is elementary: if
\((\operatorname{Ad}(\bar r(s))-1)Y\in b_v\)
for every local \(s\), induction on lower
entries, using the separated characters,
forces \(Y\in b_v\).

At each real place the unitary-adjoint fixed
space has dimension \(n(n-1)/2\):
multiplication by the pairing matrix identifies
it with skew-symmetric matrices.
The ordinary Borel has dimension \(n(n+1)/2\).
Equations (AF9), (LS8), and (LS9) give
\(t-\dim H_C^2\ge nd+1\).
Pad a generating list of \(I_R\) by zeros.
We obtain the actual ordinary presentation
\[
 R\simeq\mathcal O'[[X_1,\ldots,X_t]]/(f_1,\ldots,f_b),
                         \qquad t=b+nd+1 .                   \tag{LS10}
\]
The dual-invariant hypothesis in (AF9) is
satisfied here.
The residual matrices are triangular with
diagonal \(\Psi^{w_i}\).
Solving a twisted invariant by entries reduces
to \(\bar\varepsilon\Psi^j=1\), \(j=0,\pm1\),
or \(\bar\varepsilon\Psi^{\pm2}\delta=1\);
the latter sign comes from opposite-corner
transpose at complex conjugation.
The scalar element (AF8) excludes every case.

Suppose now \(n=3\) and \(\mathfrak e\subseteq x\).
The sign action extends to an involution of
the same \(R\), fixing \(\mathcal O'\) and \(B\).
Use \((1\pm\sigma)/2\) on tangent variables,
minimal relations, and the lifts in (LS7).
The obstruction maps split into the even and
odd cone groups.
Write \(a,g\) for the tangent dimensions,
and \(r_e,r_o\) for the relation numbers.
Equation (LS8) gives \(r_e\le h_{C,e}^2\),
\(r_o\le h_{C,o}^2\).
The tangent comparison gives
\[
 h_{C,e}^1+1\le a+h_{C,e}^0,\qquad h_{C,o}^1\le g .
\]
The extra direction is the even matrix
\(\operatorname{diag}(1,0,-1)\).
The local flags can be chosen compatible
with parity by projecting their separated
eigenvectors.
Their upper-triangular matrix units give
\(\dim b_{v,e}\ge4,\ \dim b_{v,o}\ge2\).
At a real involution the skew-matrix
identification bounds the two fixed spaces
by \(1\) and \(2\), respectively.
Apply (AF9) in each parity:
\[
 h_{C,e}^1-h_{C,e}^2\ge h_{C,e}^0+3d,\qquad
 h_{C,o}^1-h_{C,o}^2\ge h_{C,o}^0 .
\]
Hence \(a-r_e\ge3d+1\), \(g\ge r_o\).
Pad the relation lists separately to get
\[
 R\simeq
 \mathcal O'[[x_1,\ldots,x_a,y_1,\ldots,y_g]]
   /(F_1,\ldots,F_b,G_1,\ldots,G_g),\qquad a=b+3d+1 .           \tag{LS11}
\]
Here \(F_i\) are even, \(G_i\) are odd, every
linear relation coefficient is nonunit,
all entries have their prescribed parity,
\(B\) is even, and \(\mathfrak e=(y_j)\).
Assign each odd monomial to its first
occurring odd variable.
This writes \(G=My\), with every constant
coefficient of \(M\) in
\(\mathfrak m_{\mathcal O'}\).
All fields of the source's \srcid{EPres}
object have now been constructed.

\paragraph{Sources and scope.}
The actual producers under
\srcid{Development.AutomorphicArithmetic.AutomorphicVarieties}
at \srcid{alpoge-source-7cbf7a0c} are
\srcid{FrobeniusSmoothVectors.SpaceEx},
\srcid{AlgModel}, \srcid{AlgFamily},
\srcid{AlgLimit}, \srcid{PointLoc},
\srcid{DefIso}, \srcid{LocCpl},
\srcid{LocFlat}, and \srcid{LocDim}.
The obstruction and presentation calculations
expand \srcid{HasObs.pres},
\srcid{PrSetlift}, \srcid{PrCorrect},
\srcid{PrNorm}, \srcid{PrCount},
\srcid{PrMazur}, \srcid{PrDual},
\srcid{LieTMem.ep\_tan},
\srcid{EpHasobs}, \srcid{EpMazur},
and \srcid{SumLe.epres}.
This proves these specialized producers for
the arithmetic frame already constructed;
no automorphic-support implication has been
assumed in proving them.
